\documentclass[10pt]{article}
\usepackage[T1]{fontenc}
\usepackage[utf8]{inputenc}
\usepackage[french]{babel}
\usepackage{amsmath,amssymb,mathrsfs,mathpazo}
\usepackage{geometry}
\geometry{a4paper,margin=2cm}
\usepackage{graphicx,array,tabularx,tabularray,colortbl,multirow,multicol,booktabs,enumitem,pifont,ulem,fancybox,diagbox,cancel}
\usepackage{pstricks,pstricks-add,pst-plot,pst-tree,pst-node,pst-text,pst-func,pst-math,pst-bspline,pst-3dplot,pst-eucl,pst-solides3d}
\let\pscalculate\undefined
\usepackage{pst-bezier}
\usepackage{tikz,pgfplots,tkz-tab}
\usetikzlibrary{arrows,arrows.meta,patterns,calc,positioning,trees,decorations.markings,shapes.geometric,shapes.misc,shapes.arrows,intersections,decorations.pathmorphing,decorations.pathreplacing,backgrounds,shadows,babel}
\pgfplotsset{compat=1.18}
\usepgfplotslibrary{fillbetween}
\usepackage[np]{numprint}
\usepackage{listings,makecell,pas-tableur,twoopt}
\let\barre\undefined
\usepackage[pstricks]{tablvar}
\definecolor{darkgreen}{RGB}{0,100,0}
\providecommand{\red}{\color{red}}
\providecommand{\blue}{\color{blue}}
\providecommand{\green}{\color{green}}
\providecommand{\black}{\color{black}}
\providecommand{\white}{\color{white}}
\providecommand{\gray}{\color{gray}}
\providecommand{\cyan}{\color{cyan}}
\providecommand{\magenta}{\color{magenta}}
\providecommand{\yellow}{\color{yellow}}
\providecommand{\violet}{\color{violet}}
\providecommand{\orange}{\color{orange}}
\providecommand{\euro}{\texteuro}
\providecommand{\EUR}{\texteuro}
\providecommand{\ieme}{\textsuperscript{e}}
\providecommand{\degres}{\ensuremath{^\circ}}
\providecommand{\degree}{\ensuremath{^\circ}}
\providecommand{\vv}[1]{\overrightarrow{#1}}
\providecommand{\Coord}[1]{#1}
\providecommand{\decosix}{}
\providecommand{\starredbullet}{\ensuremath{\bullet}}
\pagestyle{empty}
\setlength{\parindent}{0pt}
\usepackage{adjustbox}\DeclareUnicodeCharacter{2236}{:}\DeclareUnicodeCharacter{2190}{\ensuremath{\leftarrow}}
\providecommand{\R}{\mathbb{R}}
\providecommand{\N}{\mathbb{N}}
\providecommand{\Z}{\mathbb{Z}}
\providecommand{\Q}{\mathbb{Q}}
\providecommand{\D}{\mathbb{D}}
\providecommand{\C}{\mathbb{C}}
\providecommand{\e}{{\mathrm{e}}}
\providecommand{\vect}[1]{\overrightarrow{#1}}
\providecommand{\vectt}[1]{\overrightarrow{\mathrm{#1}}}
\providecommand{\vvt}[1]{\overrightarrow{\mathrm{#1}}}
\providecommand{\np}[2][]{{#2~#1}}
\providecommand{\vv}[1]{{\overrightarrow{#1}}}
\renewcommand{\d}{{\mathrm{d}}}
\providecommand{\degres}{{^\circ}}
\providecommand{\degre}{{^\circ}}
\providecommand{\blue}{\color{blue}}
\providecommand{\red}{\color{red}}
\providecommand{\enskip}{\quad}
\providecommand{\llbracket}{\lbrack\!\lbrack}
\providecommand{\rrbracket}{\rbrack\!\rbrack}
\providecommand{\euro}{{\text{€}}}
\providecommand{\no}{{\text{n°}}}
\providecommand{\interval}[2]{{[#1;#2]}}
\providecommand{\intervFO}[2]{{[#1;#2[}}
\providecommand{\dotfill}{\ldots}

\providecommand{\ds}{\displaystyle}
\providecommand{\textding}[1]{#1}
\providecommand{\Oij}{\(\mathrm{O};\vec{\imath},\vec{\jmath}\)}
\providecommand{\Oijk}{\(\mathrm{O};\vec{\imath},\vec{\jmath},\vec{k}\)}
\providecommand{\Ouv}{\(\mathrm{O};\vec{u},\vec{v}\)}
\providecommand{\intervFF}[2]{\ensuremath{[#1;#2]}}
\renewcommand{\euro}{\ifmmode\text{\texteuro}\else\texteuro\fi}
\DeclareUnicodeCharacter{00B0}{\ifmmode^{\circ}\else\textdegree\fi}
\DeclareUnicodeCharacter{20AC}{\ifmmode\text{\texteuro}\else\texteuro\fi}
% Énoncés : Éducation nationale. Transcription et corrigés : APMEP, Denis Vergès et contributeurs. Découpage et classement : CoopMaths, Sébastien Lozano, Éric Elter et contributeurs. Adaptation de la présentation : Maths Explorer.
% CC BY-SA : https://coopmaths.fr/www/about/
% https://coopmaths.fr/alea/static/bac/2022/tex/bac_2022_05_sujet1_metropole_2.tex
\begin{document}
\noindent\textit{APMEP / CoopMaths. CC BY-SA. Présentation adaptée par Maths Explorer.}\par\medskip
\medskip

Dans l'espace rapporté à un repère orthonormé \Oijk, on considère:

$\bullet~~$ le point A de coordonnées $(-1~;~1~;~3)$,

$\bullet~~$la droite $\mathcal{D}$ dont une représentation paramétrique est: $\left\{\begin{array}{l c l}
x&=&1+2t\\y &=& 2 - t,\\z&=& 2+2t
\end{array}\right.  t \in \R$.

On admet que le point A n'appartient pas à la droite $\mathcal{D}$.

\medskip

\begin{enumerate}
\item 
	\begin{enumerate}
		\item Donner les coordonnées d'un vecteur directeur $\vect{u}$ de la droite $\mathcal{D}$.
		

		\item Montrer que le point B$(-1~;~3~;~0)$ appartient à la droite $\mathcal{D}$.
		\item Calculer le produit scalaire $\vect{\text{AB}} \cdot \vect{u}$.
	\end{enumerate}	
\item On note $\mathcal{P}$ le plan passant par le point A et orthogonal à la droite $\mathcal{D}$, et on appelle H le point d'intersection du plan $\mathcal{P}$ et de la droite $\mathcal{D}$. Ainsi, H est le projeté orthogonal de A sur la droite $\mathcal{D}$.
	\begin{enumerate}
		\item Montrer que le plan $\mathcal{P}$ admet pour équation cartésienne: $2x - y + 2z - 3 = 0$.
		\item En déduire que le point H a pour coordonnées $\left(\dfrac79~;~\dfrac{19}{9}~;~\dfrac{16}{9}\right)$.
		\item Calculer la longueur AH. On donnera une valeur exacte.
	\end{enumerate}	
\item Dans cette question, on se propose de retrouver les coordonnées du point H, projeté orthogonal du point A sur la droite $\mathcal{D}$, par une autre méthode.

On rappelle que le point B$(-1~;~3~;~0)$ appartient à la droite $\mathcal{D}$ et que le vecteur $\vect{u}$ est un vecteur directeur de la droite $\mathcal{D}$.
	\begin{enumerate}
		\item Justifier qu'il existe un nombre réel $k$ tel que $\vect{\text{HB}} = k\vect{u}$.
		\item Montrer que $k = \dfrac{\vect{\text{AB}} \cdot \vect{u}}{\left\|\vect{u}\right\|^2}$.
		\item Calculer la valeur du nombre réel $k$ et retrouver les coordonnées du point H.
	\end{enumerate}
\item On considère un point C appartenant au plan $\mathcal{P}$ tel que le volume du tétraèdre ABCH soit égal à $\dfrac89$.

Calculer l'aire du triangle ACH.

On rappelle que le volume d'un tétraèdre est donné par: $V = \dfrac13 \times \mathcal{B} \times h$ où $\mathcal{B}$ désigne l'aire d'une base et $h$ la hauteur relative à cette base.
\end{enumerate}

\bigskip

\end{document}
